% !TeX root = ../main.tex
\chapter{Archival OCR and VLM Pipeline Code}
\label{app:code_ocr}

This appendix provides the functional Python implementation for the archival optical character recognition (OCR) and vision-language model (VLM) extraction pipeline, specifically designed to process degraded physical municipal scans from the Davao City Sangguniang Panlungsod archives.

\section{Image Enhancement and Deskewing Preprocessing}
\label{sec:app_code_image_prep}

Listing~\ref{lst:code_ocr_preprocessing} illustrates the automated image restoration module that applies adaptive thresholding, bilateral noise filtering, and orientation correction to historical typewriter scans before passing them to the text extraction models.

\begin{lstlisting}[language=Python, caption={Archival Scan Preprocessing and Contrast Enhancement (\texttt{scripts/ocr\_pipeline.py})}, label={lst:code_ocr_preprocessing}]
import cv2
import numpy as np
from PIL import Image

def preprocess_archival_scan(image_path: str) -> np.ndarray:
    """Preprocesses degraded historical scans to eliminate bleed-through and shadow noise."""
    image = cv2.imread(image_path)
    gray = cv2.cvtColor(image, cv2.COLOR_BGR2GRAY)

    # Bilateral filter removes paper grain while preserving sharp typewriter edges
    denoised = cv2.bilateralFilter(gray, d=9, sigmaColor=75, sigmaSpace=75)

    # Adaptive Gaussian binarization dynamically handles uneven ink fading
    binary = cv2.adaptiveThreshold(
        denoised, 255, cv2.ADAPTIVE_THRESH_GAUSSIAN_C,
        cv2.THRESH_BINARY, blockSize=31, C=15
    )

    # Deskewing based on minimum area bounding box of detected text contours
    coords = np.column_stack(np.where(binary < 255))
    angle = cv2.minAreaRect(coords)[-1]
    if angle < -45:
        angle = -(90 + angle)
    else:
        angle = -angle

    (h, w) = binary.shape[:2]
    center = (w // 2, h // 2)
    m = cv2.getRotationMatrix2D(center, angle, 1.0)
    rotated = cv2.warpAffine(
        binary, m, (w, h), flags=cv2.INTER_CUBIC, borderMode=cv2.BORDER_REPLICATE
    )

    return rotated
\end{lstlisting}

\section{Local 4-Bit Vision-Language Model Extraction Engine}
\label{sec:app_code_vlm_engine}

Listing~\ref{lst:code_vlm_extractor} details the inference code executing the locally hosted 4-bit quantized vision-language model (\texttt{Qwen2.5-VL-3B-Instruct}) using HuggingFace Transformers and bitsandbytes. The system enforces zero hallucination and verbatim preservation of legal numerals and penalty terms.

\begin{lstlisting}[language=Python, caption={Local 4-Bit VLM Inference Engine for Archival Extraction}, label={lst:code_vlm_extractor}]
import torch
from transformers import AutoProcessor, Qwen2VLForConditionalGeneration
from PIL import Image

class ArchivalVLMReader:
    """Local 4-bit vision-language inference engine for historical ordinance digitization."""
    MODEL_ID = "Qwen/Qwen2.5-VL-3B-Instruct"

    def __init__(self, device="cuda" if torch.cuda.is_available() else "cpu"):
        self.device = device
        self.model = Qwen2VLForConditionalGeneration.from_pretrained(
            self.MODEL_ID,
            torch_dtype=torch.float16 if self.device == "cuda" else torch.float32,
            load_in_4bit=True if self.device == "cuda" else False,
            device_map="auto"
        )
        self.processor = AutoProcessor.from_pretrained(self.MODEL_ID)

    def extract_verbatim_text(self, image_path: str) -> str:
        """Extracts authentic legislative text with verbatim preservation of legal numerals."""
        image = Image.open(image_path).convert("RGB")
        prompt = (
            "You are an archival legal document transcriber. Transcribe all text on this page "
            "verbatim. Preserve exact legal numbers, monetary penalties, section headers, and "
            "punctuation. Do not correct typos, summarize, or omit clauses."
        )
        messages = [
            {"role": "user", "content": [
                {"type": "image", "image": image},
                {"type": "text", "text": prompt}
            ]}
        ]
        
        text_input = self.processor.apply_chat_template(messages, tokenize=False, add_generation_prompt=True)
        inputs = self.processor(text=[text_input], images=[image], padding=True, return_tensors="pt").to(self.device)

        with torch.no_grad():
            output_ids = self.model.generate(**inputs, max_new_tokens=2048, temperature=0.01)
        
        generated_ids = [out[len(inp):] for inp, out in zip(inputs.input_ids, output_ids)]
        return self.processor.batch_decode(generated_ids, skip_special_tokens=True)[0].strip()
\end{lstlisting}
